% ch05_brouwer.tex — Brouwer Fixed Point Theorem
\chapter{Brouwer Fixed Point Theorem}\label{ch:brouwer}

\section{Introduction and statement}

Stir your coffee in a cup, as vigorously as you like. When the liquid settles, at least one point ends up exactly where it started. This astonishing fact is a consequence of Brouwer's fixed point theorem (1911), one of the jewels of topology. L.~E.~J.\ Brouwer proved it using combinatorial methods (Sperner's lemma), and since then the theorem has found applications in game theory (Nash equilibrium), mathematical economics (Arrow-Debreu theorem), and nonlinear analysis.

Unlike the Banach principle, it is purely topological in nature: no metric
hypothesis on $T$ is required, only continuity suffices.

\begin{theorem}[Brouwer, 1911]\label{thm:brouwer}
Let $B^n = \{ x \in \R^n : \norm{x} \leq 1 \}$ be the closed unit ball in
$\R^n$. Every continuous map $f : B^n \to B^n$ has at least one fixed point.
\end{theorem}

More generally:

\begin{theorem}[Brouwer, general version]
Let $K$ be a nonempty compact convex subset of $\R^n$. Every continuous map
$f : K \to K$ has a fixed point.
\end{theorem}

\begin{remark}
The two versions are equivalent. Indeed, every nonempty compact convex subset of
$\R^n$ is homeomorphic to $B^m$ for some $m \leq n$ (if $K$ has dimension $m$).
\end{remark}

\section{The no-retraction lemma}

Brouwer's theorem is equivalent to the following topological result.

\begin{definition}[Retraction]
Let $A \subset X$ be a subspace. A \emph{retraction} of $X$ onto $A$ is a
continuous map $r : X \to A$ such that $r(a) = a$ for all $a \in A$.
\end{definition}

\begin{theorem}[No retraction]\label{thm:nonretraction}
There is no continuous retraction of $B^n$ onto $S^{n-1} = \partial B^n$.
\end{theorem}

\begin{proposition}
Theorem~\ref{thm:brouwer} and Theorem~\ref{thm:nonretraction} are equivalent.
\end{proposition}

\begin{proof}
\textbf{No retraction $\Rightarrow$ Brouwer.}
Suppose by contradiction that $f : B^n \to B^n$ continuous has no fixed point.
Then $f(x) \neq x$ for all $x \in B^n$. Define $r : B^n \to S^{n-1}$ by
sending $x$ to the intersection of the ray from $f(x)$ through $x$ with the
sphere $S^{n-1}$. This map is continuous (explicit formulas show this) and
$r(x) = x$ for $x \in S^{n-1}$, contradicting the no-retraction theorem.

\textbf{Brouwer $\Rightarrow$ No retraction.}
Suppose there exists a retraction $r : B^n \to S^{n-1}$. Define
$f : B^n \to B^n$ by $f(x) = -r(x)$. Then $f$ is continuous and
$f(B^n) \subset S^{n-1} \subset B^n$. If $f(x^*) = x^*$, then
$x^* \in S^{n-1}$ and $-r(x^*) = x^*$, so $-x^* = x^*$, hence $x^* = 0$,
contradicting $x^* \in S^{n-1}$.
\end{proof}

\section{Proof via Sperner's lemma}

We first present the combinatorial proof, due to Knaster, Kuratowski, and
Mazurkiewicz (1929), which uses Sperner's lemma.

\subsection{Sperner's lemma}

\begin{definition}[Triangulation and Sperner labeling]
Let $\sigma = [v_0, \ldots, v_n]$ be an $n$-simplex in $\R^n$. A
\emph{triangulation} of $\sigma$ is a simplicial complex $\mathcal{T}$ whose
union is $\sigma$. A \emph{Sperner labeling} is a map
$\ell : V(\mathcal{T}) \to \{0, 1, \ldots, n\}$ (where $V(\mathcal{T})$ is
the vertex set of $\mathcal{T}$) such that:
\begin{itemize}
    \item $\ell(v_i) = i$ for each vertex $v_i$ of $\sigma$;
    \item If $v$ lies on the face $[v_{i_1}, \ldots, v_{i_k}]$, then
          $\ell(v) \in \{i_1, \ldots, i_k\}$.
\end{itemize}
\end{definition}

\begin{lemma}[Sperner, 1928]\label{lem:sperner}
Let $\sigma$ be an $n$-simplex with a triangulation $\mathcal{T}$ and a Sperner
labeling. Then the number of subsimplices of $\mathcal{T}$ whose vertices carry
all labels $\{0, 1, \ldots, n\}$ is odd (and hence nonzero).
\end{lemma}

\begin{proof}[Proof in dimension $n = 1$]
Consider the segment $[v_0, v_1]$ subdivided into subsegments. The vertices are
labeled $0$ or $1$, with $v_0$ labeled $0$ and $v_1$ labeled $1$. Traversing the
segment from left to right, the number of label changes is odd (one goes from an
initial $0$ to a final $1$).
\end{proof}

\begin{proof}[Proof in general dimension (sketch)]
One uses a parity argument. Consider the dual graph: each $n$-simplex of
$\mathcal{T}$ is a vertex, and two vertices are connected if they share an
$(n-1)$-dimensional face carrying labels $\{0, \ldots, n-1\}$. One shows that
the vertices of odd degree are exactly the completely labeled simplices and
those with a boundary face carrying labels $\{0, \ldots, n-1\}$. By the
handshaking lemma and by induction (Sperner in dimension $n-1$), the number of
completely labeled simplices is odd.
\end{proof}

\subsection{Proof of Brouwer's theorem via Sperner}

\begin{proof}[Proof of Theorem~\ref{thm:brouwer}]
Let $\sigma = [e_0, e_1, \ldots, e_n]$ be the standard simplex in $\R^{n+1}$
(where $e_i$ are the basis vectors). Let $f : \sigma \to \sigma$ be continuous.
For each $m \geq 1$, consider the barycentric triangulation of $\sigma$ with
mesh $\leq 1/m$.

Define the labeling: for a vertex $v$ of the triangulation, write
$v = \sum \lambda_i e_i$ in barycentric coordinates. Set
$\ell(v) = \min\{i : f(v)_i \leq v_i\}$, where $f(v)_i$ and $v_i$ are the
$i$-th barycentric coordinates. Such an $i$ exists since
$\sum f(v)_i = \sum v_i = 1$ implies $f(v)_i \leq v_i$ for at least one $i$.

This labeling is a Sperner labeling (if $v$ lies on the face
$\{i : v_i = 0\}^c$, then $\ell(v) \in \{i : v_i > 0\}$).

By Sperner's lemma, there exists a completely labeled subsimplex $\sigma_m$.
Its vertices $v_m^0, \ldots, v_m^n$ satisfy $f(v_m^i)_i \leq (v_m^i)_i$ and
$\mathrm{diam}(\sigma_m) \leq 1/m$.

By compactness, extract a subsequence such that $v_m^i \to x^*$ for all $i$.
By continuity of $f$ and passing to the limit: $f(x^*)_i \leq x^*_i$ for all
$i$. Since $\sum f(x^*)_i = \sum x^*_i = 1$, we have $f(x^*)_i = x^*_i$ for
all $i$, hence $f(x^*) = x^*$.
\end{proof}

\section{Analytical proof}

We present an analytical proof of the no-retraction theorem based on
differential calculus.

\begin{proof}[Analytical proof of Theorem~\ref{thm:nonretraction}]
\textbf{Step 1: reduction to the $C^\infty$ case.}
By the Weierstrass approximation theorem, any continuous map
$r : B^n \to S^{n-1}$ with $r|_{S^{n-1}} = \mathrm{Id}$ can be uniformly
approximated by $C^\infty$ maps. By a truncation argument, we may assume $r$ is
$C^\infty$.

\textbf{Step 2: use of the change-of-variables formula.}
If $r : B^n \to S^{n-1}$ is $C^\infty$ with $r|_{S^{n-1}} = \mathrm{Id}$,
consider the differential form $\omega = \sum_{i=1}^n (-1)^{i-1} x_i \,
dx_1 \wedge \cdots \wedge \widehat{dx_i} \wedge \cdots \wedge dx_n$ on
$S^{n-1}$. Then $\int_{S^{n-1}} \omega = \mathrm{Vol}(S^{n-1}) \neq 0$.

But $r^*\omega = \omega$ on $S^{n-1}$, and by Stokes' theorem:
\[
    \int_{S^{n-1}} \omega = \int_{S^{n-1}} r^*\omega = \int_{B^n} d(r^*\omega).
\]
Now $r$ takes values in $S^{n-1}$, which has dimension $n-1$, so $r^*\omega$ is
exact on $B^n$ and $d(r^*\omega) = r^*(d\omega) = 0$ since
$d\omega = n \, dx_1 \wedge \cdots \wedge dx_n$ and $r$ maps into an
$(n-1)$-dimensional manifold. Contradiction.
\end{proof}

\section{Consequences of Brouwer's theorem}

\subsection{Fixed point theorem for simplices}

\begin{corollary}
Every continuous map from a simplex to itself has a fixed point.
\end{corollary}

\subsection{Borsuk--Ulam theorem (dimension 1)}

\begin{theorem}[Borsuk--Ulam in dimension 1]
For every continuous map $f : S^1 \to \R$, there exists $x \in S^1$ such that
$f(x) = f(-x)$.
\end{theorem}

\begin{proof}
Set $g(x) = f(x) - f(-x)$. Then $g(-x) = -g(x)$. If $g$ never vanishes, $g$
has constant sign, say $g(x_0) > 0$, but $g(-x_0) = -g(x_0) < 0$, a
contradiction by the Intermediate Value Theorem (on an arc connecting $x_0$ to
$-x_0$).
\end{proof}

\subsection{Ham Sandwich theorem}

\begin{theorem}[Ham Sandwich]
Let $A_1, \ldots, A_n$ be bounded measurable subsets of $\R^n$, each of positive
measure. There exists a hyperplane that bisects each $A_i$ into two parts of
equal measure.
\end{theorem}

\begin{remark}
This theorem is a consequence of the Borsuk--Ulam theorem.
\end{remark}

\section{Counter-examples: necessity of hypotheses}

\begin{example}[Non-convexity]
The rotation by angle $\pi/3$ on the circle $S^1$ is a continuous map from
$S^1$ to $S^1$ with no fixed point. The circle is compact but not convex.
\end{example}

\begin{example}[Non-compactness]
The translation $T(x) = x + e_1$ on $\R^n$ is continuous with no fixed point.
The space $\R^n$ is convex but not compact (nor bounded).
\end{example}

\begin{example}[Infinite dimension]
Consider the unit ball $B$ of $\ell^2$ and the shift
\[
    T(x_1, x_2, \ldots) = (\sqrt{1 - \norm{x}^2}, x_1, x_2, \ldots).
\]
One verifies that $T : B \to B$ is continuous but has no fixed point (if
$Tx = x$, then $x_1 = \sqrt{1 - \norm{x}^2}$ and $x_{n+1} = x_n$ for all $n$,
so $x = (c, c, c, \ldots)$ which is not in $\ell^2$ unless $c = 0$, but
$c = \sqrt{1 - 0} = 1$, contradiction). This shows that Brouwer's theorem does
not extend to infinite dimensions (see Schauder's theorem in Chapter~6).
\end{example}

\section{Brouwer's topological degree}

\begin{definition}[Brouwer degree]
Let $\Omega \subset \R^n$ be a bounded open set, $f : \overline{\Omega} \to \R^n$
continuous, and $p \notin f(\partial\Omega)$. The \emph{Brouwer degree}
$\deg(f, \Omega, p) \in \Z$ is the unique integer satisfying:
\begin{enumerate}
    \item \textbf{Normalization}: $\deg(\mathrm{Id}, \Omega, p) = 1$ if $p \in \Omega$;
    \item \textbf{Additivity}: if $\Omega = \Omega_1 \cup \Omega_2$ with
          $\Omega_1 \cap \Omega_2 = \emptyset$ and $p \notin f(\partial\Omega_i)$,
          then $\deg(f, \Omega, p) = \deg(f, \Omega_1, p) + \deg(f, \Omega_2, p)$;
    \item \textbf{Homotopy invariance}: if $H : [0,1] \times
          \overline{\Omega} \to \R^n$ is continuous with $p \notin H(t, \partial\Omega)$
          for all $t$, then $\deg(H(t, \cdot), \Omega, p)$ is constant in $t$.
\end{enumerate}
\end{definition}

\begin{theorem}[Brouwer via degree]
Let $f : B^n \to B^n$ be continuous. Then $\deg(\mathrm{Id} - f, B^n, 0) = 1$,
so $f$ has a fixed point.
\end{theorem}

\begin{proof}
The homotopy $H(t, x) = x - tf(x)$ satisfies $H(0, x) = x$ and
$H(1, x) = x - f(x)$. For $x \in S^{n-1}$ and $t \in [0,1]$:
$H(t, x) = 0$ would imply $x = tf(x)$, hence $1 = \norm{x} = t\norm{f(x)} \leq t$,
so $t = 1$ and $f(x) = x$, meaning $x$ is a fixed point on the boundary. Even
in this case, the degree is well defined and equals $1$ by normalization.

Hence $\deg(\mathrm{Id} - f, B^n, 0) = \deg(\mathrm{Id}, B^n, 0) = 1 \neq 0$,
and the equation $x - f(x) = 0$ has a solution.
\end{proof}

\section{Lefschetz fixed point theorem}

\begin{theorem}[Lefschetz]
Let $X$ be a compact absolute neighborhood retract (ANR) and $f : X \to X$
continuous. If the Lefschetz number
\[
    \Lambda(f) = \sum_{k=0}^{\infty} (-1)^k \mathrm{tr}(f_{*k} : H_k(X; \Q) \to H_k(X; \Q))
\]
is nonzero, then $f$ has a fixed point.
\end{theorem}

\begin{remark}
Brouwer's theorem is a special case: for $X = B^n$, all homology groups are
trivial except $H_0 \cong \Q$, and $f_{*0} = \mathrm{Id}$, so
$\Lambda(f) = 1 \neq 0$.
\end{remark}

\section{Exercises}

\begin{exercise}
Prove Brouwer's theorem directly (without Sperner's lemma) in dimension
$n = 1$, i.e., that every continuous $f : [0,1] \to [0,1]$ has a fixed point.
\end{exercise}

\begin{exercise}
Prove Sperner's lemma in dimension $n = 2$ using the parity argument (dual
graph).
\end{exercise}

\begin{exercise}
Use Brouwer's theorem to show that every $n \times n$ matrix with nonnegative
real entries whose column sums are $1$ has an eigenvector associated to the
eigenvalue $1$.
\end{exercise}

\begin{exercise}
Show that Brouwer's theorem implies the fundamental theorem of algebra: every
polynomial with complex coefficients of degree $n \geq 1$ has a root.
\textit{Hint: use the topological degree.}
\end{exercise}

\begin{exercise}
Let $A$ be a real $n \times n$ matrix. Show that there exist $\lambda \in \R$
and $x \in S^{n-1}$ such that $Ax = \lambda x$ (i.e., $A$ has a ``generalized
eigenvector'' on the sphere).
\end{exercise}

\begin{exercise}
Explicitly construct a continuous map $f : B^\infty \to B^\infty$ with no fixed
point, where $B^\infty$ is the closed unit ball of an infinite-dimensional
Hilbert space.
\end{exercise}
